% Part: first-order-logic
% Chapter: introduction
% Section: semantic-notions

\documentclass[../../../include/open-logic-section]{subfiles}

\begin{document}

\olfileid{fol}{int}{sem}

\olsection{المفاهيم الدلالية}

عند البحث في $\Sat{\OLId{latin-looped}{M}}{!A}[\OLId{latin-ordinary}{s}]$ جعلنا $\OLId{latin-ordinary}{s}$ تعطي قيمًا لجميع
ال!!{variable}s تيسيرًا للتعريف، وإن كنا لا نستعمل إلا قيم
ال!!{variable}s الواقعة في~$!A$. بل الذي يتوقف عليه الحكم حقًا
هو قيم المتغيرات \emph{الحرة} في~$!A$ وحدها، وسنقيم البرهان على
ذلك. فال!!{sentence}s، لخلوها من المتغيرات الحرة، لا أثر
لـ~$\OLId{latin-ordinary}{s}$ في إرضائها في !!a{structure} أو عدمه. وإذا كانت $!A$
!!a{sentence}، جاز أن نجعل $\Sat{\OLId{latin-looped}{M}}{!A}$ بمعنى
«$\Sat{\OLId{latin-looped}{M}}{!A}[\OLId{latin-ordinary}{s}]$ لكل~$\OLId{latin-ordinary}{s}$». وهذا يكافئ، كما سنرى،
$\Sat{\OLId{latin-looped}{M}}{!A}[\OLId{latin-ordinary}{s}]$ لإسناد واحد~$\OLId{latin-ordinary}{s}$ على الأقل. فإسنادات
ال!!{variable}s لازمة لإقامة تعريف إرضاء ال!!{formula}s؛ أما
ال!!{sentence}s فإرضاؤها مستقل عن إسنادات ال!!{variable}s.

بتعريف «$\Sat{\OLId{latin-looped}{M}}{!A}$» يتحدد معنى «الحالة» و«الصدق في» بالنسبة
إلى ال!!{sentence}s في منطق الرتبة الأولى. ومن
$\Sat{\OLId{latin-looped}{M}}{!A}$ لل!!{sentence}s نشتق المفاهيم الدلالية الأساسية:
الصلاحية والاستلزام وقابلية الإرضاء. فصلاحية الجملة، ورمزها
$\Entails !A$، أن ترضيها كل !!{structure}. ولزومها عن
مجموعة من ال!!{sentence}s، ورمزه $\OLId{greek-variable}{Gamma} \Entails !A$، أن كل
!!{structure} ترضي ال!!{sentence}s في~$\OLId{greek-variable}{Gamma}$ جميعًا ترضي
~$!A$ كذلك. وقابلية مجموعة من ال!!{sentence}s للإرضاء أن توجد
!!{structure} واحدة ترضي ال!!{sentence}s كلها في تلك المجموعة.

وال!!{formula}s معرفة استقراءً، والإرضاء معرّف بدوره على تركيب
ال!!{formula}s بالاستقراء؛ فلنا بهذا أن نستعمل الاستقراء لإثبات
خواص الدلالة وما بين مفاهيمها من علاقات. سنجمع بعض هذه الخواص
ونثبتها، لما فيها من فائدة في أنفسها، ولحاجتنا إليها خاصة في
البحث عن صلة الدلالة بأنساق ال!!{derivation}. ولا يتم ذلك حتى
نعرّف، على وجه دقيق استقرائي، مفاهيم وعمليات نحوية لم نتعرض لها
بعد.

\end{document}
